%! Unidad 6 — Campo magnético e inducción electromagnética: ejercicios finales (adaptados de la fuente, con datos cambiados y resueltos)

\section{Ejercicios}\label{sec:mag-ejercicios}

Salvo que se indique otra cosa, usa \(e = \qty{1.6e-19}{\coulomb}\), \(m_e = \qty{9.11e-31}{\kg}\), \(m_p = \qty{1.67e-27}{\kg}\),
\(\mu_0 = 4\pi\cdot 10^{-7}\ \unit{\tesla\m\per\ampere}\) y \(g = \qty{9.8}{\m\per\second\squared}\).

\subsection{Cuestiones}

\begin{exercise}
    Una partícula cargada se mueve en un campo magnético uniforme sin experimentar ninguna fuerza. ¿Cómo es su velocidad respecto del campo?
    \begin{solu}
        \textcolor{red}{Paralela (o antiparalela) al campo}: \textcolor{blue}{entonces \(\sen\varphi = 0\) y \(F = qvB\sen\varphi = 0\).}
    \end{solu}
\end{exercise}

\begin{exercise}
    ¿Verdadero o falso?: «La fuerza magnética que actúa sobre una carga en movimiento no hace trabajo».
    \begin{solu}
        \textcolor{red}{Verdadero.} \textcolor{blue}{La fuerza es siempre perpendicular a la velocidad (y al desplazamiento): solo cambia la dirección
        de la velocidad, no su módulo ni la energía cinética.}
    \end{solu}
\end{exercise}

\begin{exercise}
    ¿Puede ser nula la fuerza magnética sobre una partícula cargada que se mueve en un campo magnético no nulo? ¿Y la fuerza eléctrica
    sobre una partícula cargada en un campo eléctrico no nulo?
    \begin{solu}
        \textcolor{red}{La magnética sí}, \textcolor{blue}{si se mueve paralela al campo.} \textcolor{red}{La eléctrica no}: \(\vec F = q\vec E\)
        \textcolor{blue}{no depende de la velocidad y solo es nula si \(q = 0\).}
    \end{solu}
\end{exercise}

\begin{exercise}
    Un neutrón entra perpendicularmente en un campo magnético uniforme a gran velocidad. ¿Qué trayectoria sigue?
    \begin{solu}
        \textcolor{red}{Una línea recta, con velocidad constante}: \textcolor{blue}{no tiene carga y el campo no ejerce fuerza sobre él.}
    \end{solu}
\end{exercise}

\begin{exercise}
    Una carga entra de izquierda a derecha en un campo uniforme perpendicular al papel y entrante, y se desvía en el sentido de las
    agujas del reloj. ¿Qué signo tiene?
    \begin{solu}
        \textcolor{blue}{Moviéndose hacia la derecha, girar en sentido horario es curvarse hacia abajo. Para una carga positiva,
        \(\vec v\times\vec B = \vec{i}\times(-\vec{k}) = \vec{j}\): hacia arriba.} \textcolor{red}{La carga es negativa.}
    \end{solu}
\end{exercise}

\begin{exercise}
    Una espira circular está en un campo uniforme. ¿Con qué orientación es máximo el flujo que la atraviesa? ¿Con cuál es nulo?
    \begin{solu}
        \textcolor{red}{Máximo con la espira perpendicular al campo} (\(\vec S\) paralelo a \(\vec B\)); \textcolor{red}{nulo con la espira paralela al campo}
        (\(\vec S\perp\vec B\)).
    \end{solu}
\end{exercise}

\begin{exercise}
    Indica el sentido de la corriente inducida en una espira cuando se aleja de ella el polo sur de un imán, y cuando se le acerca el polo norte.
    \begin{solu}
        \textcolor{blue}{Al alejar el polo sur, la espira se opone atrayéndolo: la cara que mira al imán se convierte en un polo norte. Al acercar el
        polo norte, se opone repeliéndolo: también forma un polo norte.} \textcolor{red}{En los dos casos la corriente, vista desde el imán, circula en
        sentido antihorario.}
    \end{solu}
\end{exercise}

\subsection{Fuerzas sobre cargas y corrientes}

\begin{exercise}
    Un protón se mueve a \qty{2e7}{\m\per\second} en un campo de \qty{1.5}{\tesla} y sufre una fuerza de \qty{3e-12}{\newton}. ¿Qué ángulo
    forma su velocidad con el campo?
    \begin{solu}
        \(\sen\varphi = \dfrac{F}{qvB} = \dfrac{3\cdot 10^{-12}}{1{,}6\cdot 10^{-19}\cdot 2\cdot 10^7\cdot 1{,}5} = 0{,}625 \Rightarrow \textcolor{red}{\varphi = \ang{38.7}}\).
    \end{solu}
\end{exercise}

\begin{exercise}
    Un electrón entra con \qty{4e6}{\m\per\second} perpendicularmente en un campo de \qty{5e-3}{\tesla}. Calcula la fuerza, el radio de su
    trayectoria y su período.
    \begin{solu}
        \(F = evB = \textcolor{red}{\qty{3.2e-15}{\newton}}\); \(R = \dfrac{m v}{e B} = \textcolor{red}{\qty{4.55e-3}{\m}}\);
        \(T = \dfrac{2\pi m}{e B} = \textcolor{red}{\qty{7.16e-9}{\second}}\).
    \end{solu}
\end{exercise}

\begin{exercise}
    En una región hay un campo \(\vec E = 3\,\vec{i} + 5\,\vec{j}\ \unit{\newton\per\coulomb}\) y un campo \(\vec B = 0{,}5\,\vec{k}\ \unit{\tesla}\).
    Un electrón entra con \(\vec v = 1000\,\vec{i}\ \unit{\m\per\second}\). Calcula la fuerza y la aceleración.
    \begin{solu}
        \(\vec v\times\vec B = 1000\,\vec{i}\times 0{,}5\,\vec{k} = -500\,\vec{j}\); \(\vec E + \vec v\times\vec B = 3\,\vec{i} - 495\,\vec{j}\).
        \[ \vec F = -e\,(3\,\vec{i} - 495\,\vec{j}) = \textcolor{red}{(-4{,}8\cdot 10^{-19}\,\vec{i} + 7{,}92\cdot 10^{-17}\,\vec{j})\ \unit{\newton}} \]
        \(a = \dfrac{F}{m} = \textcolor{red}{\qty{8.69e13}{\m\per\second\squared}}\).
    \end{solu}
\end{exercise}

\begin{exercise}
    Un electrón se acelera desde el reposo con una diferencia de potencial de \qty{200}{\volt} y entra perpendicularmente en un campo de
    \qty{0.01}{\tesla}. Calcula su velocidad, el radio de la trayectoria y el período.
    \begin{solu}
        \(v = \sqrt{\dfrac{2eV}{m}} = \textcolor{red}{\qty{8.38e6}{\m\per\second}}\); \(R = \dfrac{mv}{eB} = \textcolor{red}{\qty{4.77}{\mm}}\);
        \(T = \dfrac{2\pi m}{eB} = \textcolor{red}{\qty{3.58e-9}{\second}}\).
    \end{solu}
\end{exercise}

\begin{exercise}
    En una región hay un campo magnético \(\vec B = 0{,}4\,\vec{k}\ \unit{\tesla}\) y uno eléctrico \(\vec E = 200\,\vec{j}\ \unit{\newton\per\coulomb}\).
    ¿Con qué velocidad debe entrar una partícula cargada para seguir una línea recta?
    \begin{solu}
        \textcolor{blue}{Las fuerzas deben anularse: \(\vec E + \vec v\times\vec B = 0\). Con \(\vec v = v\,\vec{i}\): \(v\,\vec{i}\times 0{,}4\,\vec{k} = -0{,}4v\,\vec{j}\), luego \(200 = 0{,}4v\).}
        \(\textcolor{red}{\vec v = 500\,\vec{i}\ \unit{\m\per\second}}\) \textcolor{blue}{(sea cual sea su carga: es un \emph{selector de velocidades}).}
    \end{solu}
\end{exercise}

\begin{exercise}
    Un electrón entra de izquierda a derecha con \qty{2e6}{\m\per\second} en un campo de \qty{2e-3}{\tesla} perpendicular al papel y entrante.
    \begin{enumerate}[label=\letraminuscula]
        \item Calcula la fuerza que actúa sobre él.
        \item Calcula el radio de su trayectoria. ¿Cómo sería la de un protón?
        \item ¿Qué campo eléctrico habría que añadir para que siguiera en línea recta?
    \end{enumerate}
    \begin{solu}
        \begin{enumerate}[label=\letraminuscula]
            \item \(\vec v\times\vec B = 2\cdot 10^6\,\vec{i}\times(-2\cdot 10^{-3}\,\vec{k}) = 4000\,\vec{j}\);
            \(\textcolor{red}{\vec F = -e\cdot 4000\,\vec{j} = -6{,}4\cdot 10^{-16}\,\vec{j}\ \unit{\newton}}\) \textcolor{blue}{(hacia abajo).}
            \item \(R = \dfrac{mv}{eB} = \textcolor{red}{\qty{5.69}{\mm}}\). \textcolor{blue}{El protón se curvaría hacia arriba (sentido antihorario) con un radio
            \(m_p/m_e \approx 1835\) veces mayor:} \textcolor{red}{\qty{10.4}{\m}}.
            \item \(\vec E = -\vec v\times\vec B = \textcolor{red}{-4000\,\vec{j}\ \unit{\newton\per\coulomb}}\).
        \end{enumerate}
    \end{solu}
\end{exercise}

\begin{exercise}
    Una bobina rectangular de 40 espiras de \qty{4}{\cm} \(\times\) \qty{6}{\cm} lleva una corriente de \qty{30}{\milli\ampere}. Calcula su
    momento magnético y el momento del par si se coloca paralela a un campo de \qty{0.3}{\tesla}.
    \begin{solu}
        \(m = NIS = 40\cdot 0{,}03\cdot 2{,}4\cdot 10^{-3} = \textcolor{red}{\qty{2.88e-3}{\ampere\m\squared}}\);
        \textcolor{blue}{paralela al campo, \(\vec S\perp\vec B\):} \(M = mB = \textcolor{red}{\qty{8.64e-4}{\newton\m}}\).
    \end{solu}
\end{exercise}

\begin{exercise}
    Un conductor de \qty{15}{\cm} lleva \qty{5}{\ampere} y forma \ang{35} con un campo horizontal. ¿Qué campo hace falta para que la fuerza
    sobre él sea de \qty{0.3}{\newton}?
    \begin{solu}
        \(B = \dfrac{F}{IL\sen\alpha} = \dfrac{0{,}3}{5\cdot 0{,}15\cdot\sen\ang{35}} = \textcolor{red}{\qty{0.697}{\tesla}}\).
    \end{solu}
\end{exercise}

\begin{exercise}
    Una varilla horizontal de \qty{25}{\cm} y \qty{15}{\g} puede deslizar sin rozamiento por dos raíles verticales, y por ella circula una
    corriente de \qty{2}{\ampere} hacia la derecha. ¿Qué campo magnético perpendicular al plano de los raíles la mantiene en reposo?
    \begin{solu}
        \textcolor{blue}{La fuerza magnética debe ser vertical hacia arriba e igual al peso: \(ILB = mg\).}
        \(B = \dfrac{0{,}015\cdot 9{,}8}{2\cdot 0{,}25} = \textcolor{red}{\qty{0.294}{\tesla}}\).
        \textcolor{blue}{Como \(\vec{i}\times(-\vec{k}) = \vec{j}\),} \textcolor{red}{el campo debe ser entrante} \textcolor{blue}{(con el eje \(Y\) hacia arriba y el \(Z\) hacia el lector).}
    \end{solu}
\end{exercise}

\subsection{Campos creados por corrientes}

\begin{exercise}
    Un solenoide de \qty{25}{\cm} tiene 500 espiras y lleva una corriente de \qty{2}{\ampere}. ¿Cuánto vale el campo en su interior? Una
    partícula cargada entra en el solenoide moviéndose a lo largo de su eje: ¿qué trayectoria sigue?
    \begin{solu}
        \(B = \mu_0\,\dfrac{N}{L}\,I = 4\pi\cdot 10^{-7}\,\dfrac{500}{0{,}25}\cdot 2 = \textcolor{red}{\qty{5.03e-3}{\tesla}}\).
        \textcolor{red}{Una recta}: \textcolor{blue}{se mueve paralela al campo y no sufre fuerza.}
    \end{solu}
\end{exercise}

\begin{exercise}
    Dos conductores paralelos y muy largos están separados \qty{10}{\cm} y llevan \qty{10}{\ampere} y \qty{5}{\ampere} en el mismo sentido.
    Calcula el campo en un punto de la recta perpendicular a ellos situado \qty{5}{\cm} a la izquierda del primero, y en el punto medio.
    \begin{solu}
        \textcolor{blue}{A la izquierda de los dos, los campos tienen el mismo sentido:}
        \(B = \dfrac{\mu_0}{2\pi}\left(\dfrac{10}{0{,}05} + \dfrac{5}{0{,}15}\right) = \textcolor{red}{\qty{4.67e-5}{\tesla}}\).

        \textcolor{blue}{En el punto medio tienen sentidos opuestos:} \(B = \dfrac{\mu_0}{2\pi}\,\dfrac{10 - 5}{0{,}05} = \textcolor{red}{\qty{2e-5}{\tesla}}\).
    \end{solu}
\end{exercise}

\begin{exercise}
    Dos conductores rectilíneos y paralelos, perpendiculares al papel, están separados \qty{1}{\m}. Por el primero circulan \qty{4}{\ampere}
    saliendo del papel.
    \begin{enumerate}[label=\letraminuscula]
        \item ¿Qué corriente debe circular por el segundo, y en qué sentido, para que el campo se anule en un punto entre ellos situado a
        \qty{80}{\cm} del primero?
        \item Calcula la fuerza por unidad de longitud entre los dos hilos.
    \end{enumerate}
    \begin{solu}
        \begin{enumerate}[label=\letraminuscula]
            \item \textcolor{blue}{Entre los hilos los campos son opuestos si las corrientes tienen el mismo sentido:}
            \(\dfrac{4}{0{,}8} = \dfrac{I_2}{0{,}2} \Rightarrow \textcolor{red}{I_2 = \qty{1}{\ampere}}\), \textcolor{red}{también saliente}.
            \item \(\dfrac{F}{L} = \dfrac{\mu_0 I_1 I_2}{2\pi d} = \textcolor{red}{\qty{8e-7}{\newton\per\m}}\), \textcolor{red}{de atracción}.
        \end{enumerate}
    \end{solu}
\end{exercise}

\begin{exercise}
    Dos conductores rectilíneos y paralelos, separados \qty{15}{\cm}, llevan \qty{8}{\ampere} y \qty{5}{\ampere}. Calcula la fuerza por unidad de
    longitud sobre cada uno si las corrientes tienen el mismo sentido y si tienen sentidos contrarios.
    \begin{solu}
        \(\dfrac{F}{L} = \dfrac{\mu_0 I_1 I_2}{2\pi d} = 2\cdot 10^{-7}\,\dfrac{8\cdot 5}{0{,}15} = \textcolor{red}{\qty{5.33e-5}{\newton\per\m}}\) \textcolor{blue}{sobre cada uno:}
        \textcolor{red}{de atracción si tienen el mismo sentido y de repulsión si son contrarias.}
    \end{solu}
\end{exercise}

\begin{exercise}
    Por un conductor horizontal fijo circulan \qty{80}{\ampere}. Otro conductor paralelo de \qty{1}{\m} y \qty{15}{\g}, situado encima, puede
    moverse verticalmente. ¿Qué corriente, y en qué sentido, debe llevar para quedar en equilibrio a \qty{4}{\mm} del primero?
    \begin{solu}
        \textcolor{blue}{La fuerza debe ser de repulsión (hacia arriba), así que} \textcolor{red}{la corriente debe tener sentido contrario}.
        \[ \dfrac{\mu_0 I_1 I_2 L}{2\pi d} = mg \;\Rightarrow\; I_1 = \dfrac{mg\,2\pi d}{\mu_0 I_2 L} = \dfrac{0{,}015\cdot 9{,}8\cdot 2\pi\cdot 0{,}004}{4\pi\cdot 10^{-7}\cdot 80} = \textcolor{red}{\qty{36.8}{\ampere}} \]
    \end{solu}
\end{exercise}

\subsection{Flujo e inducción}

\begin{exercise}
    Calcula el flujo a través de una espira de vector superficie \(\vec S = (4\,\vec{i} + 3\,\vec{j} - 12\,\vec{k})\ \unit{\cm\squared}\) en un campo
    \(\vec B = 0{,}2\,\vec{i}\ \unit{\tesla}\).
    \begin{solu}
        \(\Phi = \vec B\cdot\vec S = 0{,}2\cdot 4\cdot 10^{-4} = \textcolor{red}{\qty{8e-5}{\weber}}\).
    \end{solu}
\end{exercise}

\begin{exercise}
    Una espira de \qty{20}{\cm\squared} está en el plano \(XY\). Calcula el flujo que la atraviesa si \(\vec B = 0{,}3\,\vec{i} + 0{,}05\,\vec{k}\ \unit{\tesla}\).
    \begin{solu}
        \textcolor{blue}{\(\vec S = 2\cdot 10^{-3}\,\vec{k}\ \unit{\m\squared}\); solo cuenta la componente \(z\) del campo:}
        \(\Phi = 0{,}05\cdot 2\cdot 10^{-3} = \textcolor{red}{\qty{1e-4}{\weber}}\).
    \end{solu}
\end{exercise}

\begin{exercise}
    Una bobina de 200 espiras cuadradas de \qty{5}{\cm} de lado está perpendicular a un campo de \qty{0.4}{\tesla}. ¿Qué fem media se induce
    si gira \ang{90} en \qty{0.02}{\second}?
    \begin{solu}
        \textcolor{blue}{El flujo pasa de \(BS\) a 0:} \(|\varepsilon| = N\,\dfrac{BS}{\Delta t} = 200\,\dfrac{0{,}4\cdot 0{,}0025}{0{,}02} = \textcolor{red}{\qty{10}{\volt}}\).
    \end{solu}
\end{exercise}

\begin{exercise}
    Una bobina de 300 espiras de \qty{8}{\cm} de radio está perpendicular a un campo de \qty{0.6}{\tesla}. Calcula la fem media inducida si el
    campo se anula en \qty{0.15}{\second}.
    \begin{solu}
        \(|\varepsilon| = N\,\dfrac{B\pi r^2}{\Delta t} = 300\,\dfrac{0{,}6\cdot\pi\cdot 0{,}08^2}{0{,}15} = \textcolor{red}{\qty{24.1}{\volt}}\).
    \end{solu}
\end{exercise}

\begin{exercise}
    El flujo a través de cada espira de una bobina de 400 espiras pasa de \qty{3e-4}{\weber} a \qty{1e-4}{\weber} en \qty{0.05}{\second}. ¿Qué fem media se induce?
    \begin{solu}
        \(\varepsilon = -N\,\dfrac{\Delta\Phi}{\Delta t} = -400\,\dfrac{(1 - 3)\cdot 10^{-4}}{0{,}05} = \textcolor{red}{\qty{1.6}{\volt}}\).
    \end{solu}
\end{exercise}

\begin{exercise}
    Una espira cuadrada de \qty{4}{\cm} de lado entra a \qty{5}{\cm\per\second} en una región con un campo perpendicular de \qty{0.5}{\tesla}.
    ¿Qué fem se induce mientras entra? ¿Y cuando está entera dentro del campo?
    \begin{solu}
        \textcolor{blue}{Mientras entra, el área dentro del campo crece a razón de \(lv\):} \(\varepsilon = B l v = 0{,}5\cdot 0{,}04\cdot 0{,}05 = \textcolor{red}{\qty{1e-3}{\volt}}\).
        \textcolor{red}{Cuando está dentro, el flujo no cambia y \(\varepsilon = 0\).}
    \end{solu}
\end{exercise}

\begin{exercise}
    Una bobina de 20 espiras de \qty{5}{\cm} de radio está perpendicular a un campo que varía según \(B = 0{,}05t - 0{,}02t^2\) (SI). Calcula el
    flujo que atraviesa la bobina en función del tiempo y la fem inducida en \(t = \qty{3}{\second}\).
    \begin{solu}
        \(\Phi = N B S = 20\cdot\pi\cdot 0{,}05^2\,(0{,}05t - 0{,}02t^2) = \textcolor{red}{0{,}157\,(0{,}05t - 0{,}02t^2)\ \unit{\weber}}\).

        \(\varepsilon = -\dfrac{d\Phi}{dt} = -0{,}157\,(0{,}05 - 0{,}04t)\); en \(t = 3\): \(\varepsilon = \textcolor{red}{\qty{0.011}{\volt}}\).
    \end{solu}
\end{exercise}

\begin{exercise}
    El flujo a través de una espira de \qty{20}{\cm\squared} varía según \(\Phi = 4\cdot 10^{-4}\cos(100t)\) (SI). Calcula el campo, supuesto
    uniforme, y la fem inducida.
    \begin{solu}
        \(B = \dfrac{\Phi_{\text{máx}}}{S} = \dfrac{4\cdot 10^{-4}}{2\cdot 10^{-3}} = \textcolor{red}{\qty{0.2}{\tesla}}\);
        \(\varepsilon = -\dfrac{d\Phi}{dt} = \textcolor{red}{0{,}04\sen(100t)\ \unit{\volt}}\).
    \end{solu}
\end{exercise}

\begin{exercise}
    Una espira cuadrada de \qty{10}{\cm} de lado está en el plano \(XY\), en un campo que forma \ang{30} con el eje \(Z\) y cuyo módulo varía
    según \(B = 0{,}02\,(t^2 + 3t + 1)\) (SI). Calcula el flujo en función del tiempo y la fem en \(t = \qty{5}{\second}\). ¿En qué sentido circula
    la corriente?
    \begin{solu}
        \(\Phi = B S\cos\ang{30} = \textcolor{red}{1{,}73\cdot 10^{-4}\,(t^2 + 3t + 1)\ \unit{\weber}}\);
        \(\varepsilon = -\dfrac{d\Phi}{dt} = -1{,}73\cdot 10^{-4}(2t + 3)\); en \(t = 5\): \(|\varepsilon| = \textcolor{red}{\qty{2.25e-3}{\volt}}\).

        \textcolor{blue}{El flujo en el sentido \(+Z\) aumenta, así que el campo inducido apunta hacia \(-Z\):} \textcolor{red}{la corriente circula en sentido
        horario visto desde arriba (desde \(+Z\)).}
    \end{solu}
\end{exercise}

\subsection{Alternador y transformador}

\begin{exercise}
    Una bobina de 100 espiras cuadradas de \qty{4}{\cm} de lado gira en un campo de \qty{0.5}{\tesla}. ¿Con qué frecuencia debe girar para
    que la fem máxima sea de \qty{10}{\volt}?
    \begin{solu}
        \(\omega = \dfrac{\varepsilon_0}{NBS} = \dfrac{10}{100\cdot 0{,}5\cdot 0{,}0016} = \qty{125}{\radian\per\second}\);
        \(f = \dfrac{\omega}{2\pi} = \textcolor{red}{\qty{19.9}{\hertz}}\).
    \end{solu}
\end{exercise}

\begin{exercise}
    Un generador tiene una bobina de 20 espiras de \qty{0.05}{\m\squared} que gira a \qty{50}{\hertz} en un campo de \qty{0.4}{\tesla}. Calcula la
    fem máxima.
    \begin{solu}
        \(\varepsilon_0 = NBS\omega = 20\cdot 0{,}4\cdot 0{,}05\cdot 100\pi = \textcolor{red}{\qty{126}{\volt}}\).
    \end{solu}
\end{exercise}

\begin{exercise}
    Una bobina de 50 espiras de \qty{200}{\cm\squared} gira a \qty{600}{rpm} alrededor de un eje de su plano perpendicular a un campo de
    \qty{0.3}{\tesla}. Escribe la fem inducida en función del tiempo.
    \begin{solu}
        \(\omega = 600\cdot\dfrac{2\pi}{60} = 20\pi\ \unit{\radian\per\second}\); \(\varepsilon_0 = NBS\omega = 50\cdot 0{,}3\cdot 0{,}02\cdot 20\pi = \qty{18.8}{\volt}\).
        \[ \textcolor{red}{\varepsilon = 18{,}8\sen(20\pi t)\ \text{(V)}} \]
    \end{solu}
\end{exercise}

\begin{exercise}
    Se quiere obtener \qty{12}{\volt} a partir de la red de \qty{220}{\volt} con un transformador cuyo primario tiene 1000 espiras. ¿Cuántas
    espiras debe tener el secundario? Si por el secundario circulan \qty{2}{\ampere}, ¿qué intensidad circula por el primario?
    \begin{solu}
        \(N_2 = N_1\,\dfrac{V_2}{V_1} = 1000\,\dfrac{12}{220} \approx \textcolor{red}{55\ \text{espiras}}\);
        \(I_1 = I_2\,\dfrac{V_2}{V_1} = \textcolor{red}{\qty{0.109}{\ampere}}\).
    \end{solu}
\end{exercise}
